% ============================================================
%  SECTION 9: CONCLUSION
% ============================================================

\section{Conclusion}
\label{sec:conclusion}

We document that following Tennessee's HOPE Scholarship, public colleges
experienced a progressive shift in socioeconomic composition. The share
of bottom-quintile students increased by 1.4 percentage points
($p < 0.001$), mean parental income rank declined by 1.5 percentage
points, and the effect is monotonically ordered across all five income
quintiles. Lower-income students gained representation while
higher-income students lost it. Across 26 Tennessee public colleges,
this translates to roughly 270 additional low-income students and 740
fewer high-income students per birth cohort.

This compositional change reflects sorting across sectors rather than
enrollment expansion. Total enrollment at Tennessee public colleges is
unchanged post-HOPE ($p = 0.78$). The level decomposition confirms
that the number of bottom-quintile students increased while the number
of top-quintile students decreased. Top-quintile students who exit the
public sector do not appear at Tennessee private colleges, consistent
with out-of-state enrollment as the relevant margin.

We decompose this compositional effect into a permanent price
effect---the relative price reduction that generates progressive
sorting---and a cyclical income-sensitivity effect---the sensitivity
of that sorting to local economic conditions. Merging BEA county income
data with the Chetty panel, we show that the income-sensitivity effect
is compensated by tuition-scaling programs (Florida Bright Futures,
West Virginia PROMISE) but not by fixed-dollar programs (Tennessee
HOPE, Georgia HOPE). The shock-insurance analysis produces a second
five-for-five monotonic quintile gradient, estimated on entirely
different variation, consistent with the adoption-DiD gradient. Award
structure---not threshold height or raw generosity---determines whether
progressive sorting is sustained during economic downturns.

The cross-state evidence supports the model's comparative statics.
Massachusetts' Adams Scholarship independently produces a progressive
shift of similar magnitude ($+1.4$ pp). South Dakota's Opportunity
Scholarship produces a strongly regressive shift ($-1.5$ pp)---as the
model predicts under high eligibility thresholds. Florida and West
Virginia, whose programs predate the Chetty panel and are not testable
through adoption DiD, are testable through the shock-insurance design
and show the strongest buffering effects. The three-destination model,
extended with the price/income-sensitivity decomposition, organizes all
cases through a common set of design parameters.

Four limitations merit emphasis. First, the analysis relies on
ecological inference from college-level aggregates; we address this
through the level decomposition showing real headcount changes, not
denominator artifacts. Second, the bordering-state specification has
only four state clusters. We address this with exact $t$-based
inference and randomization inference following
\textcite{lee2026simple}, which confirm significance ($p = 0.005$;
RI $p = 0.008$) under procedures valid with a single treated state.
Unit-specific detrending attenuates the bottom-quintile estimate to
$+1.25$ pp ($p = 0.089$) and the parental-rank effect loses
significance ($-0.40$ pp, $p = 0.382$), bounding the treatment
effect and indicating that some of the estimated impact may reflect
unit-specific trends. Third, while the
preferred specification improves Rambachan--Roth sensitivity
substantially ($\Delta_\text{RM}$ robust at $\bar{M} \leq 1$,
$\Delta_\text{SD}$ robust at $M \leq 0.005$), the smoothness
restriction remains the binding constraint at moderate bounds. The
shock-insurance results provide independent corroboration through
different identifying variation. Fourth, while the pooled
shock-interaction estimates (tuition-scaling vs.\ fixed-dollar) are
precisely estimated, most individual-state shock-buffering coefficients
are not statistically significant (only Florida and West Virginia are
significant on $Q_1$). This reflects limited within-state variation:
each state contributes only 20--50 public colleges, and the
income-shock interaction is identified from county-level variation
across those colleges.

For policy, the cross-state patterns are consistent with four design
principles. First, set the eligibility threshold low enough that a
substantial share of lower-income students qualifies. Second,
concentrate the subsidy at public institutions to create a meaningful
sector price differential. Third, do not offset existing need-based
aid. Fourth, structure the award as a percentage of tuition rather than
a fixed dollar amount, so that progressive sorting is sustained when
family budgets tighten. The relevant counterfactual is not whether
merit aid is more progressive than need-based aid---it almost certainly
is not---but whether merit aid \emph{at the margin} improves the
composition and resilience of the institutions that serve the most
students.

Several extensions would strengthen the evidence. First, access to the
individual-level tax-linked enrollment records underlying
\textcite{chetty2020mobility}---available through Federal Statistical
Research Data Centers---would permit direct observation of which
students change enrollment decisions in response to merit aid,
resolving the ecological inference limitation. Second, race-specific
composition data, not available in the current public release, would
allow testing the model's prediction that the racial GPA gap
attenuates the progressive composition effect for Black students---a
prediction supported by our NLSY97 eligibility analysis but not yet
testable at the college level. Third, structural estimation of the
model's preference parameters, using Bayesian methods to formally
incorporate the NLSY97 eligibility rates as prior information, would
recover the price elasticity of public college enrollment by income
and gender---moving from reduced-form evidence to quantitative
predictions about counterfactual policies.
